% Tento soubor nahraďte vlastním souborem s obsahem práce.
%=========================================================================
% Autoři: Michal Bidlo, Bohuslav Křena, Jaroslav Dytrych, Petr Veigend a Adam Herout 2019

% Pro kompilaci po částech (viz projekt.tex), nutno odkomentovat a upravit
%\documentclass[../projekt.tex]{subfiles}
%\begin{document}

\chapter{Úvod}

V této práci se čtenář dozví informace o hře Futoshiki, její pravidla a různých
metodách, které slouží pro řešení této hry. Čtenář se dozví o konečných
automatech, které budou využití pro vytovření automatu řešitele hlavolamu.

\chapter{Futoshiki a algoritmy pro řešení}

V této kapitole jsou podrobně popsány pravidla hry Futoshiki společně s
algoritmy, které jsou využívány pro řešení hlavolamu. Začneme podrobným popisem
pravidel, které je nezbytné znát pro pochopení mechanismů hry a následného
pochopení řešících algoritmů.

Prvním a také základním algoritmem je algoritmus Backtracking. Ten funguje na
principu systematického prozkoumávání všech možných řešení metodou pokus-omyl.
Přestože toto řešení není nejvíce efektivní, je však klíčový pro pochopení
dalších sofistikovanějších metod řešení hry Futoshiki.

Následně se zaměříme na algoritmus Forward Checking a dvě jeho možné
implementace. Tento algoritmus vylepšuje princip systematického procházení
Backtrackingu o snižování počtu zkoušených kombinací pomocí předběžné kontroly.

Zároveň veškeré popsané algoritmu jsou i implementovány v programu, který
umožnuje vyřešit libovolnou přížku Futoshiki, což nám umožní demonstrovat
jejich účinnost v různých scénářích hry.

\section{Pravidla hry}

Futoshiki \cite{FutoshikiWiki}, také známá jako větší/menší (More or Less), je
logická hra pocházející z Japonska a je postavena na stejném principu jako
sudoku. Hlavolam se skládá ze čtvercové mřížky, typicky o velikosti 4x4, může
však být i větší. Pro úspešné vyřešení musí být každé hrací pole vyplněno
číslem od 1 do velikosti mřížky. Tato herní pole musí být vyplněna tak, aby se
v každém řádku i sloupci čísla neopakovala. Dále také musí platit nerovnosti,
které jsou umístěny mezi libovolné dvě hrací pole. Hra začíná s libovolným
množstvím vyplněných herních polí a nerovností tak, aby nebyla porušena
pravidla a aby existovalo alespoň jedno řešení. Následně hráč vyplňuje
zbývající herní pole, podmínky však již nemění a ani netvoří nové.

\section{Backtracking}

Backtracking \cite{PrednaskyIZU} je jednoduchý a často používaný algoritmus pro
nalezení řešení problémů. Funguje na principu systematického prozkoumávání
všech možných řešení metodou pokus-omyl. Začíná s jednou možností a pokračuje
do hloubky. V případě, že aktuálně prohledávaná cesta nefunguje (nesplňuje
podmínky,...), se vrací a zkouší jinou cestu. Takto se cyklus hledání cesty
opakuje, dokud není nalezena cesta splňující řešení nebo nejsou projity všechny
možné cesty.

\subsection{Backtracking v praxi}

Tento algoritmus je velmi jednoduchý na implementaci, avšak to vede k
nevýhodám. Backtracking zkouší všechny možnosti, i když je z počátku jasné, že
daná cesta nebude vést k řešení. Z důvodu tohoto problému se backtracking hodí
pro řešení menších problémů, například při mřížce Futoshiki o velikosti 4x4.
Jakmile se však velikost mřížky zvyšuje, efektivita algoritmu rapidně klesá.

\subsection{Backtracking a Futoshiki}

Pro řešení Futoshiki backtracking postupně zkouší dosazovat čísla do mřížky,
dokud neobjeví řešení, které splňuje všechna pravidla. Algoritmus postupuje
zleva doprava a zhora dolů v mřížce. Jakmile objeví prázdné herní pole, dosadí
do něj nějaké z čísel od 1 do velikosti hracího pole a ověří platnost pravidel.
Jestliže pravidla jsou splněna, algoritmus pokračuje na další pole. V opačném
případě se vrací na předchozí pole a zkouší se jiná čísla. Takto se algoritmus
opakuje, dokud nenalezne platné řešení nebo nevyzkouší všechny kombinace čísel.

\section{Forward Checking}

Forward checking \cite{PrednaskyIZU} vylepšuje princip systematického
prohledávání Backtrackingu o snižování počtu zkoušených kombinací. Na počátku
algoritmu se pro každé hrací pole vygeneruje doména, což je množina všech
hodnot, kterých může dané pole nabývat. Následně se na základě zvolené
strategie vybere jedno z polí, vyplní se jednou z hodnot v jeho doméně a upraví
se domény ostatních polí tak, aby byly vyloučeny hodnoty, které vedly k
porušení pravidel. Tento postup se opakuje, dokud se nenalezne řešení.

\subsection{Variace Forward Checkingu}

Uložení domén může probíhat několika způsoby. Mezi klasické způsoby patří
ukládání hodnot do pole, kde je však problematické vyloučení hodnoty z domény
nebo také ověření, zda se hodnota v doméně vyskytuje, z důvodu nutnosti projití
celého pole. Tento problém lze vyřešit pomocí hašovací množiny, která je
postavena na základě hašovací tabulky. Pro ověření, zda se daná hodnota v
doméně vyskytuje, není třeba procházet celé pole a při mazání se jednotlivé
prvky nepřeskupují. Jelikož se však pro každé hrací pole musí uložit doména, s
rostoucí velikostí mřížky se i navyšuje využití paměti. Z tohoto důvodu se
využívá bitmap, které vylepšují využití paměti tím, že je doména reprezentována
pouze jako jedno číslo, místo pole čísel. U bitmap jsme však omezeni velikostí
čísla.

Já jsem tento algoritmus implementoval dvěma způsoby: pomocí hašovacích množin
a pomocí bitmap. Implementace samotná je velice podobná, u hašovacích množin se
však využívají funkce nad množinou, zatímco u bitmap se využívají bitové
operace.

\subsection{Heuristiky Forward Checkingu}

Heuristika určuje strategii, kterou se vybírá hrací pole. Jsou dva typy
heuristik a to zvolení hracího pole a zvolení hodnoty z domény. Mezi heuristiky
výběru hracího pole patří: nejvíce omezená proměnná - hrací pole s nejmenší
doménou, nejvíce omezující proměnná - hrací pole s největším vlivem na omezení
zbývajících hracích polí. Mezi heuristiky výběru hodnoty z domény patří nejméně
omezující hodnota - vybere se hodnota, která nejméně ovlivňuje ostatní domény.

\subsection{Forward Checking a Futoshiki}

Forward Checking si pro řešení Futoshiki nejprve vygeneruje domény pro každé
hrací pole. Generacování jsem implementoval tak, že nejdříve zjistím, které
hodnoty byly již v každém řádku a v každém sloupci použity. Následně si
vytvořím doménu pro každé hrací pole tak, aby obsahovala všechny možné hodnoty
a odeberu hodnoty, které se vyskytly v daném řádku a sloupci. Nakonec také
projdu všechny nerovnosti a zaručím, aby hodnoty v doménách ovlivněných
nerovností nebyly hodnoty, které by porušily toto pravidlo. Tedy z pole, které
má být menší, musím odebrat všechny hodnoty, které jsou větší nebo rovny
maximální hodnotě z domény sousedícího pole. U opačné podmínky je třeba oddělat
hodnoty, které jsou menší nebo rovny minimální hodnotě z domény sousedícího
pole.

Po vygenerování se přejde do fáze zapisování hodnot. Pro zvolení pole jsem
zvolil heuristiku nejvíce omezená proměnná. Z domény tohoto pole poté vyberu
a dosadím první hodnotu a upravím domény ovlivněných hracích polí tak, aby
platila pravidla. Tento postup se opakuje, dokud není nalezeno řešení nebo
projity všechny možné cesty.

\section{Arc Consistency}

Moje implementace Forward Checkingu je vlastně AC-3...

\chapter{Existující formy konečných automatů}

\chapter{Dvojdimenzionální automaty}

\chapter{Implementace}

\chapter{Testování}

\chapter{Závěr}

%===============================================================================

% Pro kompilaci po částech (viz projekt.tex) nutno odkomentovat
%\end{document}
